\chapter{Реконструкция наблюдаемых}

Пусть $\mathcal D$ обозначает класс наблюдаемых, вычислительных результатов
или устойчивых эмпирических соотношений. Задача реконструкции состоит не в
том, чтобы подобрать к ним удобное толкование, а в том, чтобы найти
минимальную структуру $\mathfrak S$, из которой их организация следует с
заранее установленными условиями отказа.

\section{Устойчивое ядро данных}

Сначала из $\mathcal D$ выделяется подмножество $\mathcal D_{\mathrm{stab}}$.
В него входят только свойства, которые:
\begin{itemize}
 \item воспроизводятся независимыми способами вычисления;
 \item сохраняются при малых допустимых деформациях;
 \item не зависят от послепроверочного выбора нормировки;
 \item допускают явное описание погрешности или области применимости.
\end{itemize}

До выделения $\mathcal D_{\mathrm{stab}}$ построение модели запрещено: иначе
аппарат начинает объяснять шум и случайные особенности представления.

\section{Первичная атрибуция}

Для каждого $O\in\mathcal D_{\mathrm{stab}}$ строится предварительная карта
\begin{equation}
 \mathfrak A_0(O)\in
 \{\mathrm{car},\mathrm{phase},\mathrm{int},\mathrm{mix},\mathrm{trans}\}.
\end{equation}

Значение карты определяется чувствительностью наблюдаемого к контролируемым
изменениям каркасного, фазового и внутреннего секторов. Словесная аналогия не
считается достаточным основанием для атрибуции.

\section{Реконструкция минимальной структуры}

Ищется такой кандидат $\mathcal I^*$, что
\begin{equation}
 \mathcal I^*\Longrightarrow
 \mathfrak A_0(\mathcal D_{\mathrm{stab}})
\end{equation}
при минимальном числе независимых предположений. Несколько наблюдаемых должны
стягиваться к общему источнику, а не получать по отдельному правилу вывода.

\section{Проверка устойчивости}

Для допустимой деформации $\delta\in\Delta_{\mathrm{adm}}$ требуется
сохранение доминирующего секторного профиля:
\begin{equation}
 \mathfrak A_0(O)\simeq\mathfrak A_0(\delta O).
\end{equation}

Если малая деформация полностью меняет атрибуцию, исследование возвращается к
выделению устойчивого ядра данных или к выбору более точной модели.

\section{Смешанные наблюдаемые}

Для смешанного наблюдаемого недостаточно формальной суммы независимых частей.
Необходимо предъявить общий оператор или функционал, порождающий перекрёстный
вклад. Схематически
\begin{equation}
 O_{\mathrm{mix}}=O_{\mathrm{car}}+O_{\mathrm{phase}}+
 O_{\mathrm{int}}+O_{\times},
\end{equation}
где $O_{\times}$ должен быть выведен, а не подобран после сравнения с данными.

\chapter{Межъязыковая согласованность}

Одна структура должна допускать несколько согласованных редукций:
\begin{equation}
 R_{\mathrm{geom}}(\mathcal I^*)\sim
 R_{\mathrm{spec}}(\mathcal I^*)\sim
 R_{\mathrm{corr}}(\mathcal I^*).
\end{equation}

Знак $\sim$ означает совпадение инвариантного содержания, а не буквальное
тождество формул. Если три языка требуют трёх несвязанных источников, гипотеза
единой структуры не проходит проверку.

\chapter{Критерий содержательности}

Реконструкция считается нетривиальной только тогда, когда найденное ядро
ограничивает ещё не использованные наблюдаемые или запрещает часть возможных
форм смешения. Схема, допускающая любое продолжение, не обладает
предсказательной силой.

Полный ход реконструкции имеет вид
\begin{equation}
 \mathcal D\longrightarrow\mathcal D_{\mathrm{stab}}
 \longrightarrow\mathfrak A_0
 \longrightarrow\mathcal I^*
 \longrightarrow\Delta_{\mathrm{adm}}
 \longrightarrow\text{новые ограничения}.
\end{equation}

\chapter{Сравнение конкурирующих моделей}

Для моделей $M_1$ и $M_2$, объясняющих один класс данных, применяется единая
ведомость отбора.

\section{Совместимость режима}

Модель должна помещать наблюдаемое в тот структурный режим, который следует
из деформационной чувствительности. Неверная режимная локализация не может
компенсироваться лучшим численным совпадением.

\section{Структурная экономия}

Предпочтение получает модель, которая объясняет больше устойчивых свойств
меньшим числом независимых сущностей и правил. Скрытые исключения и отдельные
коэффициенты для каждого сектора считаются усложнением.

\section{Устойчивость}

Результат должен сохраняться при допустимых изменениях параметризации,
базиса, нормировки и вычислительной схемы. Неустойчивый результат получает
статус локального совпадения.

\section{Переносимость}

Сильная модель сохраняет одно инвариантное содержание в геометрическом,
спектральном и корреляционном языках. Если толкование приходится полностью
заменять при каждой смене языка, общий источник не установлен.

\section{Различающая сила}

Предпочтение получает модель, запрещающая больше неверных продолжений и
дающая новые проверки, не использованные при построении. Гибкость без
ограничений считается недостатком, а не преимуществом.

\chapter{Защита от подгонки}

До вычисления целевой величины фиксируются:
\begin{enumerate}
 \item допустимые входные данные;
 \item класс разрешённых формул и преобразований;
 \item способ нормировки;
 \item контрольные перестановки и исключения элементов;
 \item порог успеха и критерий провала.
\end{enumerate}

После получения результата запрещается расширять класс формул только для
сохранения совпадения. Новая гипотеза должна открываться как отдельная ветвь с
новым условием проверки.

Следующий сегмент должен описать допустимые деформации, устойчивые режимы и
границу между слабым изменением одной структуры и переходом к иной структуре.